\documentclass{article}

\usepackage{amsfonts}
\usepackage{booktabs}
\usepackage{pl-syntax}
\usepackage{metavarsx}

% ---------------------------------------------------------------------------
% Meta-variable declarations (metavarsx).
%
% \NewMetaVariable{<key>}[<\command>]{<display>}
%
% declares a meta-variable and a command that typesets it. Without the
% optional argument the command is \<key>. From <display> the package derives
%   set          \makefirstuc{<display>}
%   set of sets  \widehat{<set>}
%   universe     \mathbb{<set>}
% Any derived form can be replaced with \RenewMetaSet, \RenewMetaSetSet, or
% \RenewMetaUniverse.
% ---------------------------------------------------------------------------

% The simplest case: key e, command \e, displayed as e.
\NewMetaVariable{e}{e}

% \x and \n are free, so the key can double as the command name.
\NewMetaVariable{x}{x}
\NewMetaVariable{n}{n}

% \v is already the caron accent, so pick a different command name.
\NewMetaVariable{val}[\val]{v}

% Greek letters have no capital to derive, so set the forms by hand. The
% universe and set of sets are rebuilt from the new set automatically.
\NewMetaVariable{ty}[\ty]{\tau}
\RenewMetaSet{ty}{T}

\NewMetaVariable{ctx}[\ctx]{\Gamma}
\RenewMetaSet{ctx}{\mathcal{G}}
\RenewMetaUniverse{ctx}{\mathbf{Ctx}}

% Multi-letter names: \makefirstuc capitalizes inside \mathit, giving Msg,
% but \mathbb{\mathit{Msg}} would print as plain italic, so replace the
% universe.
\NewMetaVariable{msg}[\msg]{\mathit{msg}}
\RenewMetaUniverse{msg}{\mathbf{Msg}}


\title{Example Usage of the \texttt{pl-syntax} and \texttt{metavarsx} Packages}
\author{Andrew Hirsch \and Josh Acay}


\begin{document}
\maketitle

\section{The Simply Typed Lambda Calculus}

Read the comments in the \texttt{.tex} source for an explanation of how
the following is generated.

\begin{figure}[h]
  \begin{syntax}
    % Define a syntactic category with no right-hand side.
    \abstractCategory[Variables]{x, y, z}

    % Define a syntactic category from a set.
    \categoryFromSet[Naturals]{n}{\mathbb{N}}

    % Start a category definition. Optional input is the category
    % description, the other declares the meta variable(s) you will
    % use to refer to this category.
    \category[Expressions]{e}
    % Each alternative is introduced on its own line.
    \alternative{x}
    \alternative{n}
    \alternative{\lambda x . e}
    \alternative{e_1 e_2}

    \category[Values]{v}
    \alternative{\lambda x . e}

    % You can add vertical spacing between categories to visually group them.
    \separate
    % You can pass the amount of space explicitly if you want to manually control it:
    % \separate[5ex]

    \category[Types]{\tau}
    \alternative{\mathbf{Nat}}
    % You can place alternatives on a new line.
    \\
    \alternative{\tau_1 \rightarrow \tau_2}

    \category[Contexts]{\Gamma}
    \alternative{x_1 : \tau_1, \ldots, x_n : \tau_n}
    % You can use words to describe more complicated properties of
    % definitions if you don't want to use BNF.
    \note{no repeats}
    \note{unordered}


    % Alternative layout if you want more space.
    \category[Expressions]{e}
    \groupleft{
      \alternative{x}
      \alternative{n}
    }
    \groupleft{
      \alternative{\lambda x . e}
      \alternative{e_1 e_2}
    }
  \end{syntax}

  \caption{Terms and Types of the Simply Typed Lambda Calculus}
  \label{fig:abstract-syntax}
\end{figure}

It is possible to group multiple definition on one line:
\begin{syntax}
  \groupleft{
    \abstractCategory[Variables]{x, y, z}
    \hfill
    \quad
    \categoryFromSet[Naturals]{n}{\mathbb{N}}
    \quad
    \hfill
    \category[Expressions]{e}
    \alternative{x}
    \alternative{n}
    \alternative{\lambda x . e}
    \alternative{e_1 e_2}
  }

\end{syntax}

\section{Meta-Variables with \texttt{metavarsx}}

A meta-variable such as $e$ comes with related notation: a set of
expressions $E$, the universe of all expressions $\mathbb{E}$, and
sometimes a set of such sets. The \texttt{metavarsx} package declares
all of these from one definition, so they stay consistent across a
document and can be restyled in one place.

\subsection{Declaring a Meta-Variable}

Declarations go in the preamble:
\begin{verbatim}
\NewMetaVariable{e}{e}            % key e, command \e
\NewMetaVariable{val}[\val]{v}    % key val, command \val
\end{verbatim}
The first argument is a key that names the variable inside the package.
The optional second argument is the command to define; without it the
command is \verb|\|\textit{key}. The last argument is how the variable
itself is displayed.

\subsection{The Four Forms}

Suffixes on the command select the form:

\begin{center}
  \begin{tabular}{@{}lll@{}}
    \toprule
    Input           & Output         & Meaning \\
    \midrule
    \verb|\e|       & \e             & one expression \\
    \verb|\e*|      & \e*            & a set of expressions \\
    \verb|\e**|     & \e**           & a set of sets of expressions \\
    \verb|\e!|      & \e!            & all expressions \\
    \verb|\showcasex{\e}| & \showcasex{\e} & all three at once \\
    \bottomrule
  \end{tabular}
\end{center}

Every form works in both text and math mode, and a space after it in
running text is kept: ``each \e* is finite'' is written
\verb|each \e* is finite|.

The suffixes must follow the command directly. \verb|\e *| is the
variable $e$ followed by a literal asterisk. A \verb|!| takes precedence
over stars, so \verb|\e!*| gives the universe and discards the star.

\subsection{Overriding the Derived Forms}

By default the set is the display with its first letter capitalized, the
universe is the set in blackboard bold, and the set of sets is the set
with a wide hat. These defaults do not suit every variable, and three
commands replace them:
\begin{verbatim}
\RenewMetaSet{key}{...}
\RenewMetaSetSet{key}{...}
\RenewMetaUniverse{key}{...}
\end{verbatim}
They take the key, not the command. The default universe and set of
sets are computed from the set each time they are used, so replacing
the set alone updates both.

\paragraph{Greek letters.} \verb|\makefirstuc| has nothing to capitalize
in \verb|\tau|, so every form would print as $\tau$. After
\verb|\RenewMetaSet{ty}{T}| the forms are \showcasex{\ty}, and the set
of sets is \ty**.

\paragraph{Multi-letter names.} For \verb|\NewMetaVariable{msg}[\msg]{\mathit{msg}}|
the set capitalizes inside \verb|\mathit| as expected, but
\verb|\mathbb{\mathit{Msg}}| prints as plain italic, and
\texttt{amsfonts} provides blackboard bold for capital letters only.
Replacing the universe gives \showcasex{\msg}.

\paragraph{Fully custom forms.} The context variable replaces both the
set and the universe: \showcasex{\ctx}.

\subsection{Choosing Command Names}

Without the optional argument, the declaration defines the command
\verb|\|\textit{key}, and it fails if that command already exists. Text accents and symbols already use many one-letter
commands, among them
\begin{center}
  \verb|\b  \c  \d  \H  \i  \j  \k  \l  \o  \P  \r  \S  \t  \u  \v|
\end{center}
Give these variables an explicit command, as with \verb|\val| for $v$
above.

\subsection{Meta-Variables in a Syntax Definition}

The two packages combine directly. Figure~\ref{fig:metavars-syntax}
repeats the grammar of Figure~\ref{fig:abstract-syntax} with the
meta-variable commands, so changing a declaration in the preamble
changes the figure and the prose together.

\begin{figure}[h]
  \begin{syntax}
    \abstractCategory[Variables]{\x, y, z}
    \categoryFromSet[Naturals]{\n}{\mathbb{N}}

    \category[Expressions]{\e \in \e!}
    \alternative{\x}
    \alternative{\n}
    \alternative{\lambda \x . \e}
    \alternative{\e_1\, \e_2}

    \category[Values]{\val \in \val!}
    \alternative{\lambda \x . \e}

    \separate

    \category[Types]{\ty \in \ty!}
    \alternative{\mathbf{Nat}}
    \alternative{\ty_1 \rightarrow \ty_2}

    \category[Contexts]{\ctx \in \ctx!}
    \alternative{\x_1 : \ty_1, \ldots, \x_k : \ty_k}
    \note{no repeats}
    \note{unordered}
  \end{syntax}

  \caption{Figure~\ref{fig:abstract-syntax} written with \texttt{metavarsx}}
  \label{fig:metavars-syntax}
\end{figure}

Prose uses the same commands. For example: a context \ctx assigns types
to variables; a set of values $\val* \subseteq \val!$ is closed if every
\val in it has no free variables; and a typing environment for a module
is a set of contexts \ctx* drawn from \ctx!.

\end{document}
